\section{Crypto Agility}


\begin{definition}[Crypto Agility]
A system is \emph{crypto-agile} if it can add, replace, or deprecate cryptographic algorithms without requiring a coordinated, simultaneous upgrade of all participants.
\end{definition}

Crypto agility in a blockchain context requires:

\begin{enumerate}[leftmargin=2em]
\item \textbf{Algorithm negotiation}: Signatures, key exchanges, and hashes are identified by algorithm IDs, not hardcoded.
\item \textbf{Multi-algorithm support}: The verification path can check multiple signature types per transaction/block.
\item \textbf{Graceful deprecation}: Old algorithms can be disabled via governance without hard fork.
\item \textbf{Key rotation}: Validators can rotate to new key types without losing their stake or identity.
\end{enumerate}

\subsection{Algorithm Registry}

The Lux P-Chain maintains a registry of supported algorithms:

\begin{center}
\begin{tabular}{llll}
\toprule
\textbf{ID} & \textbf{Algorithm} & \textbf{Purpose} & \textbf{Status} \\
\midrule
\texttt{0x01} & secp256k1 (ECDSA) & EVM transactions, addresses & Active \\
\texttt{0x02} & BLS12-381 & Consensus aggregation & Active \\
\texttt{0x03} & ML-DSA-65 (FIPS 204) & PQ identity signatures & Active \\
\texttt{0x04} & ML-DSA-87 (FIPS 204) & PQ high-security signatures & Reserved \\
\texttt{0x05} & SLH-DSA-128s (FIPS 205) & Hash-based fallback & Reserved \\
\texttt{0x06} & Pulsar (Module-LWE) & PQ anonymous threshold & Active \\
\texttt{0x07} & ML-KEM-768 (FIPS 203) & PQ key encapsulation & Active \\
\bottomrule
\end{tabular}
\end{center}

New algorithms are added by P-Chain governance proposal. Deprecation follows the same process with a mandatory grace period.

